\chapter{Weil Sheaves, Weights, and Determinant Weights}
\label{chap:weights}

The purpose of this chapter is to introduce the language in which the
Riemann hypothesis becomes stable under geometric operations.  Two
distinctions must remain visible.  A Weil sheaf carries an action of the
discrete Weil group, whereas a sheaf descended to a finite field carries
continuous descent under a profinite Galois group.  Also, purity relative
to one field isomorphism with \(\mathbf C\) is weaker than the
all-conjugates assertion eventually required for integral Frobenius
polynomials.

Throughout \(k_0=\Fq\), \(k=\Fbar\), and
\[
  X=X_0\times_{k_0}k
\]
for a separated finite-type \(k_0\)-scheme \(X_0\).  We use geometric
Frobenius; it acts on \(\Ql(1)\) by \(q^{-1}\).

\section{Scheme Frobenius, Galois Frobenius, and Weil data}

Let
\[
  \Fr_{X_0}^{(q)}:X_0\longrightarrow X_0
\]
be the \(q\)-power scheme Frobenius: it is the identity on the
underlying topological space and sends a local function \(f\) to
\(f^q\).  Because the ground field is \(\Fq\), this is an
\(\Fq\)-morphism.  Define
\[
  F_X
  =
  \Fr_{X_0}^{(q)}\times_{\Spec\Fq}\id_{\Spec k}
  :X\longrightarrow X.
\tag{8.1}\label{eq:scheme-frobenius}
\]
This is a \(k\)-morphism and a universal homeomorphism; when \(X_0\) is
smooth its differential is zero.  It need not be an isomorphism of
schemes, but pullback by it is an equivalence of the \'etale topos.

This scheme map must not be confused with the semilinear Galois
automorphism.  Let
\[
  \gamma:a\longmapsto a^{1/q}
\]
be geometric Galois Frobenius and let
\(\Gamma_X=\id_{X_0}\times\gamma\).  If
\(\Fr_X^{\mathrm{abs}}\) denotes absolute \(q\)-power Frobenius on
\(X\), then
\[
  F_X=\Fr_X^{\mathrm{abs}}\circ\Gamma_X.
\tag{8.2}\label{eq:two-frobenius-comparison}
\]
Absolute Frobenius induces the identity equivalence on the \'etale
topos.  Consequently a datum
\(F_X^*\cF\simeq\cF\) is canonically the same, on the \'etale topos, as
a datum \(\Gamma_X^*\cF\simeq\cF\).  Under Galois descent the latter is
the action of \emph{geometric} Galois Frobenius.  This comparison is
also why the endomorphism induced by the scheme map \(F_X\) on
cohomology is the geometric-Galois Frobenius action used below.

\begin{definition}
\label{def:weil-sheaf}
A \emph{constructible Weil sheaf} on \(X_0\) is a pair
\((\cF,\varphi)\), where \(\cF\) is a constructible
\(\Qbarl\)-sheaf on \(X\) and
\[
  \varphi:F_X^*\cF\xrightarrow{\sim}\cF.
\]
A Weil complex is defined in the same way in the constructible derived
category.
\end{definition}

Iteration defines, at every degree \(n\) closed point \(x_0\), a
conjugacy class
\[
 F_{x_0}:\cF_{\bar x}\longrightarrow\cF_{\bar x}.
\]
Via \eqref{eq:two-frobenius-comparison}, this is geometric Galois
Frobenius.  On the constant sheaf it is the identity; on
\(\Qbarl(1)\) it is \(q^{-n}\), and on \(\Qbarl(-1)\) it is \(q^n\).
Thus \(\Qbarl(m)\) has weight \(-2m\).

If \(X_0\) is connected, its Weil group fits into
\[
 1\to\pi_1(X,\bar x)\to W(X_0,\bar x)
 \xrightarrow{\deg}\mathbf Z\to0,
\]
with geometric Frobenius of degree \(1\).  Lisse Weil sheaves are
continuous finite-dimensional representations of this group, with the
usual \(\ell\)-adic topology on the geometric subgroup and the discrete
topology on \(\mathbf Z\).

\begin{warning}
\label{warn:weil-versus-etale}
A Weil sheaf need not descend to \(X_0\).  Descent requires the action
to extend continuously from \(\mathbf Z\) to
\(\widehat{\mathbf Z}=\Gal(k/k_0)\).  The Fourier argument uses scalar
Weil twists which need not meet that condition, so the trace formula
and the six operations must be available for Weil objects themselves.
\end{warning}

\section{Punctual purity and mixedness}

Fix \(\iota:\Qbarl\xrightarrow{\sim}\mathbf C\).

\begin{definition}
\label{def:iota-purity}
Let \(\cF_0\) be a constructible Weil sheaf and \(w\in\mathbf R\).
\begin{enumerate}
\item It is \emph{punctually \(\iota\)-pure of weight \(w\)} if every
local eigenvalue \(\alpha\) at a degree \(n\) closed point satisfies
\[
 |\iota(\alpha)|=q^{nw/2}.
\]
\item It is \emph{\(\iota\)-mixed} if it has a finite filtration whose
quotients are punctually pure.
\item It has weights at most \(w\) if all quotient weights are at most
\(w\), and weights at least \(w\) analogously.
\end{enumerate}
\end{definition}

For a complex \(\cK_0\), the derived convention is
\[
 \cK_0\text{ has weights }\le w
 \Longleftrightarrow
 \cH^j(\cK_0)\text{ has weights }\le w+j.
\]

\begin{definition}
\label{def:all-embedding-purity}
An algebraic number \(\alpha\in\Qbar\) is a \(q^n\)-Weil number of
weight \(w\) if
\[
 |\sigma(\alpha)|=q^{nw/2}
\]
for every \(\sigma:\Qbar\hookrightarrow\mathbf C\).  A sheaf is pure
in the all-embeddings sense if all local eigenvalues are algebraic and
satisfy this condition.
\end{definition}

\begin{proposition}[Elementary permanence]
\label{prop:elementary-weight-permanence}
Upper bounds are preserved by subobjects, quotients, extensions, finite
sums, and finite Frobenius-stable filtrations.  Tensor products add pure
weights, duals negate them, a twist by \(m\) subtracts \(2m\), and
pullback preserves punctual weights.
\end{proposition}

\begin{proof}
Stalks are exact and characteristic polynomials multiply in short exact
sequences.  Tensor eigenvalues are products, dual eigenvalues are
inverses, and a twist by \(m\) multiplies a degree \(n\) eigenvalue by
\(q^{-mn}\).  Pullback raises Frobenius to the same power by which it
multiplies the residue degree.
\end{proof}

\section{Rank one: the independent base case}

\begin{foundation}
\label{found:rank-one-purity}
\textbf{Rank-one purity capsule.}
Let \(X_0\) be normal and geometrically connected, and let \(\cL_0\) be
a rank-one lisse Weil sheaf.  For every \(\iota\), it is punctually pure
of some real weight.

The associated character
\(\chi:W(X_0,\bar x)\to\Qbarl^\times\) has finite image on
\(\pi_1(X,\bar x)\) by the rank-one Katz--Lang theorem.  This reduces,
through a sufficiently ample normal curve, to geometric class field
theory and Lang's theorem for a generalized Picard variety; it uses no
higher-rank weight result.  If \(s\) has degree one and \(a=\chi(s)\),
then at a degree \(n\) point
\[
 \chi(F_{x_0})=\zeta_{x_0}a^n
\]
with \(\zeta_{x_0}\) a root of unity.  Hence
\[
 |\iota\chi(F_{x_0})|
 =q^{nw_\iota/2},
 \qquad
 w_\iota=2\log_q|\iota(a)|.
\]
See \cite[Section~1.3]{DeligneWeilII},
\cite[Chapter~I]{KiehlWeissauer}, and
\cite{StanfordL19,StanfordL22}.
\end{foundation}

\section{Determinant weights}

\begin{definition}
\label{def:determinant-weight}
For an irreducible lisse Weil sheaf \(\cF_0\) of rank \(r\), define
\[
 d_\iota(\cF_0)=\frac1r\,w_\iota(\det\cF_0).
\]
For a geometrically semisimple sheaf, take the determinant weights of
its irreducible constituents with multiplicity.
\end{definition}

\begin{lemma}[Determinant-weight arithmetic]
\label{lem:det-weight-arithmetic}
Duality negates determinant weights; a rank-one twist of weight \(c\)
adds \(c\); the largest weight of \(\cF^{\otimes m}\) is \(m\) times
that of \(\cF\); and exterior-power weights are the corresponding sums,
subject to constituent ranks.
\end{lemma}

\begin{proof}
After semisimplification, tensor and exterior constructions are
representations of the product of general linear groups on the
isotypic blocks.  Their central characters are the indicated monomials;
rank-one purity changes products into sums of weights.
\end{proof}

\section{Arithmetic and geometric monodromy}

Let
\[
 G_{\mathrm{geom}}
 =\overline{\rho(\pi_1(X,\bar x))}^{\,\mathrm{Zar}}
 \subseteq\operatorname{GL}(V)
\]
and let \(G_{\mathrm{arith}}\) be generated by this group and a
positive-degree Frobenius element.

\begin{proposition}[Central element of positive degree]
\label{prop:positive-degree-center}
After a connected finite \'etale cover and a finite extension of
\(k_0\), \(H=G_{\mathrm{geom}}^\circ\) is semisimple.  Moreover
\(Z(G_{\mathrm{arith}})(\Qbarl)\) contains an element of positive
degree, and
\[
 Z(G_{\mathrm{arith}})(\Qbarl)\xrightarrow{\deg}\mathbf Z
\]
has finite kernel and cokernel.
\end{proposition}

\begin{proof}
Geometric semisimplicity first makes \(H\) reductive: its unipotent
radical has a fixed vector on every irreducible constituent and
therefore acts trivially; faithfulness kills it.

After taking a power of a positive-degree element \(s\), its action on
\(H\) is inner.  Choose \(g\in H\) inducing that action and put
\(z=g^{-1}s\); then \(z\) centralizes \(H\).  If \(H\) had a nontrivial
algebraic character, the splitting generated by \(H\) and \(z\) would
extend it to a rank-one character of a finite-index Weil group with
infinite geometric image, contradicting
\cref{found:rank-one-purity}.  Thus the connected reductive group \(H\)
has no characters, hence is semisimple and has finite center.

Raise \(z\) so that it acts trivially on the finite component group of
\(G_{\mathrm{geom}}\).  For \(u\in G_{\mathrm{geom}}\),
\[
 c_u=zuz^{-1}u^{-1}\in H.
\]
Since \(z\) centralizes \(H\), direct conjugation shows
\(c_u\in Z(H)\).  This finite-valued cocycle factors through the
component group.  Raising \(z\) to the exponent of \(Z(H)\) kills all
\(c_u\), so a positive power of \(z\) is central in
\(G_{\mathrm{arith}}\).  The kernel of degree is contained in the
finite group \(Z(G_{\mathrm{geom}})\), and the positive degree of \(z\)
makes the cokernel finite.  See
\cite[Chapter~I, Section~3]{KiehlWeissauer} and
\cite{StanfordL22}.
\end{proof}

\begin{corollary}[Dense pullback preserves determinant weights]
\label{cor:dense-pullback-det-weights}
Let \(f:Y_0\to X_0\) have dense image between normal geometrically
irreducible schemes, and let \(\cF_0\) be geometrically semisimple and
lisse.  Then \(\cF_0\) and \(f^*\cF_0\) have the same determinant
weights.
\end{corollary}

\begin{proof}
Choose compatible positive-degree central elements in the two
monodromy groups and replace them by powers until degrees and images
agree.  They act by the same scalar on a constituent and on every
summand after pullback.  The logarithmic absolute value of that scalar,
divided by the degree and normalized by rank, is the determinant
weight.
\end{proof}

\section{The dependency boundary}

\begin{formalization}
\label{form:weights-interface}
\textbf{Reusable interface.}
\begin{enumerate}
\item The \(q\)-power scheme Frobenius \(F_X\), separately from the
semilinear Galois automorphism \(\Gamma_X\), together with
\eqref{eq:two-frobenius-comparison}.
\item Weil objects, effective descent, fixed-\(\iota\) purity, and
all-embeddings purity as distinct notions.
\item Elementary permanence, rank-one purity, determinant weights, and
dense-pullback invariance.
\item The semisimple finite-center proof of a positive-degree central
element.
\end{enumerate}
No result here uses \(Rf_!\)-preservation of weights or smooth-proper
purity.
\end{formalization}
